\clearpage\begingroup\SourceLeaf{92b}\HistoricalHead{87}{R}{GAUSS AN ENCKE.}
\begin{center}[11.]\end{center}
\begin{center}\textsc{Gauss} an \textsc{Encke}.\quad Göttingen, 14. Juni 1834.\end{center}
\noindent\hspace*{\parindent}Für die gewogentliche Mittheilung der Berliner magnetischen Beobachtungen
vom 21. und 22. März habe ich Ihnen, mein theuerster Freund, noch
meinen verbindlichsten Dank abzustatten. Man erkennt darin, wie Sie auch
selbst schon bemerkt haben, die Spur der in den Frühstunden des zweiten
Tages hier wahrgenommenen Anomalien, aber auch zugleich die Notwendigkeit
\textls[120]{viel} engerer Zwischenzeiten, wenn man dergleichen Anomalien und ihren
Parallelismus an verschiedenen Orten vollkommen übersehen will. Die Zahlen
der hiesigen Beobachtungen würden Ihnen gern zu Dienst stehen; allein
ihre öffentliche Bekanntmachung würde wohl nur \textls[120]{in dem} Fall von Interesse
sein können, wenn in irgend einem dritten Orte viel detaillirtere Beobachtungen
gemacht sind, als in Berlin.

Am 4. und 5. Mai sind hier die Beobachtungen noch viel detaillirter gemacht
als das erste mahl, nämlich durchweg von 10 zu 10 Minuten, ja wenigstens
zur Hälfte noch enger, nämlich von 5 zu 5 Min. und bei besonderen
Vorkommnissen anhaltend von Minute zu Minute. Beim Beobachten von 10
zu \(10'\) würden also die Aufzeichnungen so stehen
\clearpage\SourceLeaf{93}\HistoricalHead{88}{L}{PHYSIK. BRIEFWECHSEL.}
\newcommand{\LTimePair}[6]{\begin{array}{@{}r@{\;}r@{\;}r@{}}#1&#2&#3\\#4&#5&#6\end{array}}
\begin{center}\small
\(\text{Schema A.}\quad
\left.\begin{array}{@{}r@{\;}l@{\quad}l@{}}
\left.\LTimePair{8^{\mathrm{h}}}{57'}{49''}{}{58}{10}\right\}&\text{Mittel}&\text{. . .}\quad\text{gilt für }8^{\mathrm{h}}\,58'\\[.7ex]
\left.\LTimePair{}{58}{49}{}{59}{10}\right\}&\text{Mittel}&\text{. . . . . .}\quad8\;59\\[.7ex]
\left.\LTimePair{}{59}{49}{9}{0}{10}\right\}&\text{Mittel}&\text{. . . . . .}\quad9\;0\\[.7ex]
\left.\LTimePair{}{0}{49}{}{1}{10}\right\}&\text{Mittel}&\text{. . . . . .}\quad9\;1\\[.7ex]
\left.\LTimePair{}{1}{49}{}{2}{10}\right\}&\text{Mittel}&\text{. . . . . .}\quad9\;2\\[-.5ex]
\cline{1-1}
9\quad2\;49\quad&&\\
\text{etc. etc.}\quad&&
\end{array}\right\}\quad
\begin{gathered}\text{Mittel aus allen}\\\text{als für }9^{\mathrm{h}}\,0'\text{ gültig}\\\text{betrachtet.}\end{gathered}\)
\end{center}
\noindent\hspace*{\parindent}Bei Aufzeichnungen von 5 zu 5 Min. ist wohl mit halben Minuten beobachtet,
namentlich von H[errn] \textsc{Listing} und H[errn] Prof. \textsc{Weber}. Z. B.
\begin{center}\small
\(\text{Schema B.}\quad
\begin{array}{@{}r@{\quad}c@{\quad}r@{}}
\LTimePair{10^{\mathrm{h}}}{58'}{49''}{}{59}{10}&\text{. . . . . . . .}&10^{\mathrm{h}}\;59'\;0''\\[-.5ex]\cline{1-1}
\LTimePair{}{59}{19}{}{59}{40}&\text{. . . . . . . .}&10\;59\;30\\[-.5ex]\cline{1-1}
\LTimePair{}{59}{49}{11}{0}{10}&\text{. . . . . . . .}&11\;0\;0\\[-.5ex]\cline{1-1}
\LTimePair{}{0}{19}{}{0}{40}&\text{. . . . . . . .}&11\;0\;30\\[-.5ex]\cline{1-1}
\LTimePair{}{0}{49}{}{1}{10}&\text{. . . . . . . .}&11\;1\;0
\end{array}\)
\end{center}
\noindent Ich selbst habe gewöhnlich nur von 10 zu \(10'\) nach A, zuweilen von 5 zu \(5'\)
theils nach A theils nach B beobachtet.

In diesen Tagen, besonders am zweiten (5. Mai), haben sich nun merkwürdige
Anomalien gezeigt, von denen man bei Aufzeichnungen von Stunde
zu Stunde kaum eine Idee bekommen haben würde. Ich muss es daher als
etwas sehr interessantes betrachten, dass diesmal correspondirende Beobachtungen
mit beinahe eben so viel Detail da sind. H[err] \textsc{Sartorius} hat nemlich
in Waltershausen unter Beistand von zwei Gehülfen an einem zwar kleineren
\clearpage\SourceLeaf{94}\HistoricalHead{89}{R}{GAUSS AN ENCKE.}\HistoricalFoot{XI1.}{12}
\noindent aber dem hiesigen ähnlichen Apparat beobachtet, und die Mittel von 10 zu
\(10'\), wie sie aus dem Schema A sich ergeben, mir mitgetheilt. Der Beobachtungsort
zwischen Meinungen und Münnerstadt liegt ein paar Zeit-Minuten
östlich von Göttingen.

H[err] \textsc{Sartorius} hat in den ersten Stunden versäumt, wieder nach der
Marke zu sehen, welche er dann um \(10^{\mathrm{h}}\,40'\) zu seinem Schrecken bedeutend
verstellt fand. Von da an hat er jede Stunde mehreremahle nachgesehen und
nöthigenfalls nachgeholfen. Die ersten Beobachtungen sind mithin zu cassieren;
ich habe sie jedoch punctirt mitgezeichnet.

Sie werden nun gewiss mit freudiger Ueberraschung die graphische Darstellung
der Göttinger und Waltershauser Beobachtungen (jene schwarz, diese
roth) nach gleichem Maassstabe gezeichnet in dem anliegenden Blatt[\textsuperscript{*})],\begingroup\renewcommand{\thefootnote}{}
\footnotetext{[\textsuperscript{*}) Das Blatt wurde an A. \textsc{von Humboldt} weitergegeben, vergl. S. 90 unten.]}
\endgroup\ untere
Hälfte, betrachten. Die grossen schnell wechselnden Anomalien erkennt man
alle wieder und selbst an den meisten kleinern findet man die kenntliche Spur.
Ja ich halte für wahrscheinlich, dass die Harmonie noch frappanter ausfallen
würde, wenn die Waltershauser Beobachtungen in den kritischen Stunden
auch anstatt von 10 zu 10 Min. (wo sich kleinere Krausheiten verwischen)
von 5 zu 5 Minuten gemacht wären. Ich habe H[errn] \textsc{Sartorius} aufgefordert
mir noch die Resultate für jede einzelne Minute, die vorhanden sind, anstatt
der extrahirten Mittel, zu schicken, und hoffe, dass sich dann auch die Spur
der \textls[120]{hier mit entschiedenster Gewissheit} den 5. Mai Morgens \(8^{\mathrm{h}}\text{—}8^{\mathrm{h}}\,20'\)
bemerkten Einsenkung finden wird. Ich bemerke nur noch, dass die Göttinger
Zahlen \textls[120]{nicht} nach Mitteln aus je 5 Werthen, sondern direkt nach den \textls[120]{einzelnen}
vorhandenen Minuten zuerst in eine Zeichnung mit 5mahl grösserem
Maassstabe, und danach nach dem Augenmaasse verjüngt (in die correspondirenden
Quadrate) eingetragen sind.

Zur Vergleichung habe ich in der oberen Hälfte des Blattes nun auch
noch die hiesigen Beobachtungen vom 21. 22. März nach \textls[120]{demselben} Maassstabe
gezeichnet und darunter die mir von Ihnen mitgetheilten Berliner Zahlen,
ohne jedoch letztere Punkte durch Linien zu verbinden, wobei zu viel Willkührlichkeit
sich hätte einmischen müssen.

Ich glaube, dass durch die erwähnten Thatsachen sich ganz unzweifelbar
\clearpage\SourceLeaf{95}\HistoricalHead{90}{L}{PHYSIK. BRIEFWECHSEL.}
\noindent herausstellt, dass nicht bloss die grösseren Ausschweifungen, sondern selbst die
kleineren und in kurzen Fristen wechselnden Fluctuationen \textls[120]{nicht} local, sondern
durch Kräfte, welche bis in sehr grosse Entfernungen hinaus wirken,
hervorgebracht sind, und das Interesse für die gleichzeitig an vielen Orten
anzustellenden Beobachtungen wird mithin umsomehr gesteigert. Am 21. u.
22. Juni soll hier wieder fleissig beobachtet werden, obwohl ich dabei den
Abgang des H[errn] Dr. \textsc{Weber} schmerzlich vermissen werde. H[err] \textsc{Listing}
wird wahrscheinlich diesmal noch theilnehmen können; auch ein anderer Zuhörer
von mir hat sich bisher einzuüben versucht. H[err] Hofr. \textsc{Harding} hat
solches aber noch zu sehr versäumt, um diesmal schon auf seinen Beistand
zählen zu können.

Unsre grosse galvanische Kette ist seit 6 Wochen wieder ganz in Ordnung,
\textsc{Schumacher} hat sie noch spielen sehen. Eine daneben gehängte Glasglocke
tönen zu machen, einen daneben gestellten auf etwas schmaler Basis
stehenden Körper umzuwerfen, einen Wecker auszulösen u. dergl., ist ganz
leicht[\textsuperscript{*})].\begingroup\renewcommand{\thefootnote}{}
\footnotetext{[\textsuperscript{*}) Dieser Satz steht, mit roter Tinte geschrieben, am Rande des Briefes.]}
\endgroup\ Der Multiplicator hat 200 Umwindungen, und die ganze Drahtlänge
der Kette wird gegen 9000 pariser Fuss sein. Dennoch ist die Wirkung
noch zu stark für alphabetische Telegraphie, selbst wenn man nur Ein Plattenpaar
wie 1 pr[eussischer] Thaler gross und \textls[120]{reines Brunnenwasser} nimmt,
was anfangs Bewunderung erregt, aber doch ganz in der Ordnung ist.

Eine recht accurate Wiederholung der Intensitätsbestimmung hoffe ich
nun bald vornehmen zu können. Ich muss dann nur einmal anhaltend 4—5 Tage
ganz dabei bleiben können, und leider ist meine Zeit durch mancherlei
Geschäfte sehr zersplittert.

Die Stäbe, in Uslar bestellt, sind nunmehr alle fertig, obwohl noch nicht
hier und noch nicht gehärtet. Die 8 grossen (jeder à 25 {\fontspec{FreeSerif.otf}℔}) sind aber hier
und schon gehärtet und gestrichen. Zusammengesetzt werden sie ein schönes
Streichmittel bilden.

Herrn \textsc{von Humboldt} bitte ich unter meiner gehorsamsten Empfehlung
die graphische Darstellung mitzutheilen, wobei ich nur zu entschuldigen bitte,
dass sie nicht schön ist. Ich hatte sie zunächst bloss für meinen eigenen Gebrauch
gezeichnet, glaubte aber, dass es Ihnen interessant sein würde, davon
\clearpage\SourceLeaf{96}\HistoricalHead{91}{R}{GAUSS AN ENCKE.}\HistoricalFoot{}{12\textsuperscript{*}}
\noindent Kenntniss zu nehmen, und die Zeit fehlte, um eine saubere Zeichnung anzufertigen.
Finden Sie es angemessen, so können Sie auch der K. Akademie
von dem Inhalt dieses Briefes, was Sie wollen, communiciren.

\textsc{Schumacher} zeigt mir an, dass \textsc{Oerstedt} die Absicht habe, im Laufe des
Sommers, vielleicht in Gesellschaft von \textsc{Hansteen} auch hieher zu kommen.
\textsc{Schumacher} hat hier für sich einstweilen einen Apparat wie der kleinere
(mit 1 {\fontspec{FreeSerif.otf}℔}diger Nadel) bestellt; ein anderer von \textls[120]{derselben} Dimension wie
der hiesige grosse ist vor einigen Tagen für Leipzig fertig geworden, und bereits
dahin abgeschickt. Wenn ich H[errn] \textsc{Weber} recht verstanden habe, ist
auch einer für Dresden bestellt.

Ich muss in Beziehung auf die Zeichnungen noch bemerken, dass sie
überall nur nach der nominellen Zeit (Uhrzeit) gemacht sind. In Göttingen
war dies auf ein Paar Secunden genau mittlere Zeit; in Waltershausen war
die Uhr etwa 1—2 Minuten voraus, was mit dem Meridianunterschied zusammen
eine Verschiebung von etwa 3—4 Minuten nöthig machen würde, um die
dortigen Beobachtungen den Göttinger gleichzeitig zu machen. Mir deucht,
dass man dieses selbst in der Zeichnung an den hervortretenden Stellen wieder
erkennt. Ich möchte behaupten, dass
\begin{list}{}{\setlength{\leftmargin}{\parindent}\setlength{\rightmargin}{0pt}\setlength{\topsep}{0pt}\setlength{\partopsep}{0pt}}
\item[] wenn in Waltershausen an einem eben so feinen Apparat beobachtet
wäre, wie der hiesige ist (an dem dortigen entspricht 1 Scalentheil oder
Millimeter \(52''\,5\), also \(\frac{5}{2}\)mal so viel wie hier), und die Aufzeichnungen in
den letzten 6 Stunden von Minute zu Minute gemacht wären, man daraus
den Meridianunterschied bis auf einen \textls[120]{kleinen Bruchtheil einer Minute}
würde ableiten können, vorausgesetzt, dass die dabei thätigen
Kräfte den hiebei in Frage kommenden Raum in unmerklicher Zeit
durchdringen.
\end{list}
\noindent\hspace*{\parindent}Ob diese Voraussetzung wahr ist oder nicht, ist freilich noch unbekannt;
allein ich lebe der Hoffnung, dass wenn vielleicht in nicht langer Zeit ähnliche
Etablissements wie das hiesige in Altona, Copenhagen, Christiania, Berlin
und Russland (vielleicht auch Nord-Amerika) bestehen, wir über diese transcendent-interessante
Frage bald ins Reine kommen werden.

Erhalten Sie Ihr freundschaftliches Andenken
\begin{flushright}Ihrem treu ergebenen\end{flushright}
\noindent\hspace*{\parindent}Göttingen, 14. Juni 1834.\hfill C. F. \textsc{Gauss}.\hspace*{3em}
\clearpage\SourceLeaf{97a}\HistoricalHead{92}{L}{PHYSIK. BRIEFWECHSEL.}
\noindent\hspace*{\parindent}P. S. Damit Sie mich nicht missverstehen, muss ich noch eine Bemerkung
beifügen. Ich erwarte gar nicht, dass die Waltershauser Beobachtungen,
wären sie ebenso detaillirt wie die hiesigen, genau die Copie von letzteren
darbieten würden, sondern finde vielmehr sehr natürlich, wenn die Wirkungen,
je entfernter vom eigentlichen Heerd (den wir doch der Wahrscheinlichkeit
nach weit im Norden von Göttingen zu suchen haben) sich desto mehr verflachen.
Die Zukunft wird uns darüber schon Licht geben.
\clearpage\endgroup
